\documentclass[12pt, a4paper]{article}
\usepackage[utf8]{inputenc}
\usepackage[english]{babel}
\usepackage[margin=2cm]{geometry}
\usepackage{graphicx}
\usepackage{float} %use the option [H]
\graphicspath{ {images/} }
\usepackage{amsthm} %lets us use \begin{proof}
\usepackage{amsmath}
\DeclareMathOperator{\arcsinh}{arcsinh}
\usepackage{amssymb} %gives us the character
\usepackage{CJKutf8}
\usepackage[export]{adjustbox}
\setlength{\parindent}{0cm} % starting spaces
\setlength{\parskip}{1em} % paragraph width
\usepackage{array}
\usepackage{tabularx}
\title{\textbf{Calculus HW6}}
\author{王弘禹}
\date\today
\begin{document}
\begin{CJK*}{UTF8}{bsmi}
\linespread{1.5}
\maketitle
\newcommand{\st}[1]{\section*{#1}}
\newcommand{\sst}[1]{\subsection*{#1}}
\newcommand{\ssst}[1]{\subsubsection*{#1}}
\newcommand{\dsp}{\displaystyle}
\newcommand{\dx}{\,dx}
\newcommand{\dphi}{\,d\phi}
\newcommand{\dtheta}{\,d\theta}
\newcommand{\dy}{\,dy}
\newcommand{\dz}{\,dz}
\newcommand{\dr}{\,dr}
\st{15.3}
\sst{15.}
\ssst{\textbf{Problem Statement}}
Evaluate the given integral by changing to polar coordinates.
\begin{align}
    &\iint_R \arctan{(y/x)}\,dA,\\
    &\text{where } R =\{(x,y)\mid 1\le x^2+y^2\le4,\,0\le y \le x\} 
\end{align}
\ssst{\textbf{Solve.}}
Let $x=r\cos \theta$, $y=r\sin \theta \implies \arctan{\dfrac{y}{x}}=\theta$
Because $0\le y\le x$, we know the area is in first quadrant and under the line $x=y$.
\begin{align}
    &\iint_R \arctan{(y/x)}\,dA = \int_1^2\int_0^{\frac{\pi}{4}}r\theta \dtheta dr=\dfrac{3}{64}\pi^2
\end{align}
\sst{19.}
\ssst{\textbf{Problem Statement}}
Use a double integral to find the area of the region D.
\ssst{\textbf{Solve.}}
The intersect of two circles is $(\sin {(\pi/4)},\pi/4)$ and $(\sin \pi,\pi)$
\begin{align}
    \text{area of D } = &2 \int_{\pi/4}^{\pi} \int_0^{\sin\theta}r\dr d\theta\\
    =& \int_{\pi/4}^{\pi}\sin^2\theta \dtheta\\
    =&-\sin\theta \cos \theta\Big|_{\frac{\pi}{4}}^{\pi}+\int_{\pi/4}^{\pi}(1-\sin^2 \theta)\dtheta\\
    \implies \text{area of D} =&\dfrac{1}{2}(\dfrac{3}{4}\pi+\dfrac{1}{2})=\dfrac{3}{8}\pi+\dfrac{1}{4}
\end{align}

\sst{41.}
\ssst{\textbf{Problem Statement}}
Evaluate the iterated integral by converting to polar coordinates.
\begin{equation}
    \int_0^{1/2}\int_{\sqrt{3}y}^{\sqrt{1-y^2}}xy^2\dx dy
\end{equation}
\ssst{\textbf{Solve.}}
\begin{align}
    &x=\sqrt{1-y^2 } \implies x^2+y^2=1 \implies \text{a circle}\\
    &x=\sqrt{3}y \implies \text{a line with }m=\tan\theta=1/\sqrt{3} \implies \theta = \pi/6
\end{align}
Fortunately, $y=\dfrac{1}{2}=1\sin{(\dfrac{\pi}{6})}$. The area iterated integral is:
\begin{align}
    &\int_0^{1/2}\int_{\sqrt{3}y}^{\sqrt{1-y^2}}xy^2\dx dy\\
    =&\int_0^{\pi/6}\int_0^1 r^4\cos\theta \sin^2\theta\dr d\theta\\
    =&\dfrac{1}{5}\int_0^{\pi/6} \sin^2\theta \,d(\sin \theta)=\dfrac{1}{15\cdot 8}=\dfrac{1}{120}
\end{align}

\sst{49.}
\ssst{\textbf{Problem Statement}}
Use polar coordinates to combine the sums
\begin{equation}
    \int_{1/\sqrt{2}}^{1}\int_{\sqrt{1-x^2}}^x xy\dy dx + \int_1^{\sqrt{2}}\int_0^x xy\dy dx+\int_{\sqrt{2}}^{2}\int_0^{\sqrt{4-x^2}}xy\dy dx
\end{equation}
into one double integral. Then evaluate the double integral.
\ssst{\textbf{Solve.}}
\begin{align}
    &y=\sqrt{1-x^2} \implies x^2+y^2 = 1 \implies \text{a circle with radius 1}\\
    &y=x \implies \text{a line with }\tan \theta = 1 \\
    &y=0 \implies \text{x-axis}\\
    &y=\sqrt{4-x^2} \implies x^2+y^2 = 4 \implies \text{a circle with radius 2}
\end{align}
So the sums can be represented as this one double integral:
\begin{equation}
    \int_0^{\pi/4}\int_1^2 r^3\cos\theta\sin\theta \dr d\theta=\dfrac{15}{4}\int_0^{\pi/4}\sin\theta \,d{\sin\theta}=\dfrac{15}{16}
\end{equation}

\end{CJK*}
\end{document}

